%% ============================================================
%% Vorwissenschaftliche Arbeit - Hauptdatei
%% BG Zehnergasse, Wiener Neustadt
%% Vorlage mit APA6-Zitierstil
%% ============================================================
%%
%% Optionen fuer \documentclass:
%%   german/english  - Sprache (Standard: german)
%%   oneside/twoside - Ein-/Zweiseitig (Standard: oneside)
%%   color/black     - Farbig/Schwarz-Weiss (Standard: color)
%%
\documentclass[german,oneside,color]{bgzvwa}

% === Literatur-Datei ===
\addbibresource{literatur.bib}

% === Glossar vorbereiten ===
\makeglossaries

% === Glossar- und Abkuerzungseintraege ===
%% ============================================================
%% Glossar- und Abkuerzungseintraege
%% Verwendung im Text: \gls{vwa}, \gls{apa}, \gls{bmbwf}
%% ============================================================

% --- Abkuerzungen (acronym) ---
\newacronym{vwa}{VWA}{Vorwissenschaftliche Arbeit}
\newacronym{apa}{APA}{American Psychological Association}
\newacronym{bmbwf}{BMBWF}{Bundesministerium f{\"u}r Bildung, Wissenschaft und Forschung}

% --- Glossareintraege ---
\newglossaryentry{forschungsfrage}{
  name={Forschungsfrage},
  description={Die zentrale Fragestellung einer wissenschaftlichen Arbeit,
    die durch systematische Untersuchung beantwortet werden soll.
    Eine gute Forschungsfrage ist klar formuliert, abgrenzbar und
    innerhalb des gegebenen Rahmens beantwortbar}
}

\newglossaryentry{primaerquelle}{
  name={Prim{\"a}rquelle},
  description={Originaldokument oder Originaldaten, die als unmittelbarer
    Beleg f{\"u}r eine wissenschaftliche Aussage dienen, im Gegensatz
    zu Sekund{\"a}rquellen, die bereits eine Interpretation enthalten}
}

\newglossaryentry{abstract}{
  name={Abstract},
  description={Kurze Zusammenfassung einer wissenschaftlichen Arbeit,
    die Thema, Methode, Ergebnisse und Schlussfolgerungen in
    komprimierter Form wiedergibt}
}


% ============================================================
% Metadaten der VWA
% ============================================================
\title{Titel der Vorwissenschaftlichen Arbeit}
\verfasserIn{Vorname Nachname}
\klasse{8A}
\betreuungsperson{Mag.\ Vorname Nachname}
\fachbereich{Geschichte und Politische Bildung}
\abgabejahr{2025}

% Optional: Schuldaten aendern (Standardwerte sind bereits gesetzt)
% \schule{BG Zehnergasse}
% \schuladresse{Zehnergasse 15, 2700 Wiener Neustadt}
% \schullogo{bgz-logo.png}

% ============================================================
\begin{document}

% === Frontmatter (roemische Seitenzahlen) ===
\frontmatter

% === Titelseite ===
\maketitle

% === Inhaltsverzeichnis ===
\tableofcontents

% === Selbststaendigkeitserklaerung ===
\selbststaendigkeitserklaerung

% === Kurzfassung (Deutsch) ===
%% ============================================================
%% Kurzfassung (Deutsch)
%% Umfang: 1000-1500 Zeichen (ca. 150-250 Woerter)
%% ============================================================

\chapter*{Kurzfassung}
\addcontentsline{toc}{chapter}{Kurzfassung}

% Hinweis: Die Kurzfassung soll folgende Punkte abdecken:
% - Thema und Forschungsfrage
% - Methodik
% - Zentrale Ergebnisse
% - Schlussfolgerung

Die vorliegende Vorwissenschaftliche Arbeit befasst sich mit dem Thema
\glqq [Titel der Arbeit]\grqq{} und untersucht die zentrale Forschungsfrage,
[Forschungsfrage hier einfuegen]. Das Ziel der Arbeit besteht darin,
[Ziel der Arbeit beschreiben].

Zur Beantwortung der Forschungsfrage wurde eine [Methodik beschreiben,
\zB Literaturanalyse, qualitative Interviews, vergleichende Analyse]
durchgefuehrt. Die theoretischen Grundlagen stuetzen sich auf
[wesentliche Quellen oder Theorien nennen].

Die Ergebnisse zeigen, dass [zentrale Ergebnisse zusammenfassen].
Insbesondere konnte festgestellt werden, dass [wichtigste Erkenntnis].

Zusammenfassend laesst sich sagen, dass [Schlussfolgerung formulieren].
Die Arbeit leistet damit einen Beitrag zum Verstaendnis von
[Themenbereich] und zeigt moegliche Ansatzpunkte fuer weitergehende
Untersuchungen auf.

\vskip 5mm
\noindent\textbf{Schlagwoerter:} Schlagwort 1, Schlagwort 2, Schlagwort 3, Schlagwort 4


% === Abstract (Englisch) ===
%% ============================================================
%% Abstract (English)
%% Length: 1000-1500 characters (approx. 150-250 words)
%% ============================================================

\chapter*{Abstract}
\addcontentsline{toc}{chapter}{Abstract}

% Note: The abstract should cover:
% - Topic and research question
% - Methodology
% - Key findings
% - Conclusion

This pre-scientific paper examines the topic of \glqq [Title of the paper]\grqq{}
and investigates the central research question of [insert research question].
The aim of this paper is to [describe the objective].

To answer the research question, a [describe methodology, e.g., literature
analysis, qualitative interviews, comparative analysis] was conducted.
The theoretical framework is based on [mention key sources or theories].

The findings indicate that [summarize central results]. In particular,
it was determined that [most important finding].

In conclusion, [formulate conclusion]. This paper thus contributes to
the understanding of [topic area] and identifies potential avenues
for further research.

\vskip 5mm
\noindent\textbf{Keywords:} keyword 1, keyword 2, keyword 3, keyword 4


% ============================================================
% === Hauptteil (arabische Seitenzahlen) ===
% ============================================================
\mainmatter

%% ============================================================
%% Einleitung
%% ============================================================

\chapter{Einleitung}
\label{chap:einleitung}

% Hinweis: Die Einleitung fuehrt in das Thema ein und umfasst
% typischerweise 2-3 Seiten. Sie enthaelt Hinfuehrung zum Thema,
% Forschungsfrage, Methodik und Aufbau der Arbeit.

[Hinfuehrung zum Thema: Warum ist das Thema relevant? Welche aktuelle
Bedeutung hat es? Hier wird das Interesse der Leserin / des Lesers geweckt
und der Kontext der Arbeit hergestellt.]

[Persoenliche Motivation oder gesellschaftliche Relevanz des Themas
erlaeutern.]


\section{Forschungsfrage}
\label{sec:forschungsfrage}

Die zentrale Forschungsfrage dieser Vorwissenschaftlichen Arbeit lautet:

\begin{quote}
\textit{[Hier die Forschungsfrage als vollstaendigen Fragesatz formulieren,
\zB \glqq Inwiefern beeinflusste die Industrialisierung die soziale Struktur
in Wiener Neustadt im 19.~Jahrhundert?\grqq]}
\end{quote}

% Optional: Untergeordnete Fragestellungen
Zur Beantwortung dieser Fragestellung werden folgende Teilfragen untersucht:
\begin{itemize}
  \item [Teilfrage 1]
  \item [Teilfrage 2]
  \item [Teilfrage 3]
\end{itemize}


\section{Methodik}
\label{sec:methodik}

Zur Beantwortung der Forschungsfrage wird in dieser Arbeit
[Methodik beschreiben] eingesetzt. [Erklaeren, warum diese Methode
gewaehlt wurde und wie sie angewendet wird.]

% Beispiele fuer Methoden:
% - Literaturanalyse: Auswertung von Fachliteratur und Primaerquellen
% - Qualitative Interviews: Leitfadengestuetzte Interviews mit ExpertInnen
% - Vergleichende Analyse: Gegenueberstellung von Fallbeispielen
% - Empirische Erhebung: Fragebogen, Beobachtung, Experiment

[Konkrete Beschreibung der Vorgehensweise, verwendete Quellen und
Materialien.]


\section{Aufbau der Arbeit}
\label{sec:aufbau}

Die vorliegende Arbeit gliedert sich in [Anzahl] Kapitel.

Nach dieser Einleitung werden in Kapitel~\ref{chap:theoretische-grundlagen}
die theoretischen Grundlagen dargestellt. Kapitel~\ref{chap:analyse}
widmet sich [Beschreibung]. In Kapitel~\ref{chap:ergebnisse} werden
die Ergebnisse praesentiert. Abschliessend fasst Kapitel~\ref{chap:fazit}
die zentralen Erkenntnisse zusammen und beantwortet die Forschungsfrage.

%% ============================================================
%% Hauptteil
%% Der Hauptteil kann in mehrere Kapitel aufgeteilt werden.
%% Diese Datei enthaelt exemplarisch drei Kapitel:
%%   - Theoretische Grundlagen
%%   - Analyse / Untersuchung
%%   - Ergebnisse
%%
%% Bei Bedarf koennen die Kapitel auch in eigene Dateien ausgelagert
%% und einzeln per \input{} eingebunden werden.
%% ============================================================


% ============================================================
\chapter{Theoretische Grundlagen}
\label{chap:theoretische-grundlagen}
% ============================================================

% Hinweis: Dieses Kapitel stellt den wissenschaftlichen Hintergrund dar.
% Hier werden relevante Theorien, Begriffe und der Stand der Forschung
% erlaeutert. Alle verwendeten Quellen muessen korrekt zitiert werden.

[Einfuehrung in den theoretischen Rahmen der Arbeit. Welche Begriffe
und Konzepte sind zentral?]


\section{Begriffsbestimmung}
\label{sec:begriffe}

[Definition und Abgrenzung der wichtigsten Begriffe, die in der Arbeit
verwendet werden.]

% Beispiel fuer ein indirektes Zitat (Paraphrase) im APA-Stil:
% Laut \textcite{nachname2023} ist [Inhalt der Paraphrase].
% Oder: [Inhalt der Paraphrase] \parencite{nachname2023}.


\section{Stand der Forschung}
\label{sec:forschungsstand}

[Darstellung des aktuellen Forschungsstandes zum Thema. Welche
Arbeiten und Erkenntnisse gibt es bereits? Wo bestehen Forschungsluecken?]

% Beispiel fuer ein direktes Zitat im APA-Stil:
% \textcite[S.~42]{nachname2023} stellt fest: \glqq [Zitat]\grqq.


\section{Theoretischer Rahmen}
\label{sec:theorie}

[Vorstellung der theoretischen Grundlage, auf der die Arbeit aufbaut.
Welche Modelle oder Theorien werden herangezogen?]


% ============================================================
\chapter{Analyse}
\label{chap:analyse}
% ============================================================

% Hinweis: In diesem Kapitel wird die eigentliche Untersuchung
% durchgefuehrt. Hier werden die Methoden angewendet und die
% gesammelten Daten oder Quellen ausgewertet.

[Einfuehrung in das Analysekapitel. Was wird untersucht und wie?]


\section{Materialauswahl und Vorgehensweise}
\label{sec:material}

[Beschreibung der verwendeten Materialien (Quellen, Daten, Texte)
und der konkreten Vorgehensweise bei der Analyse.]


\section{Durchfuehrung der Analyse}
\label{sec:durchfuehrung}

[Darstellung der eigentlichen Analyse. Hier werden Argumente
entwickelt, Quellen interpretiert oder Daten ausgewertet.]

% Beispiel fuer eine Abbildung:
% \begin{figure}[htbp]
%   \centering
%   \includegraphics[width=0.8\textwidth]{beispiel-grafik}
%   \caption{Beschreibung der Abbildung}
%   \label{fig:beispiel}
% \end{figure}

% Beispiel fuer eine Tabelle:
% \begin{table}[htbp]
%   \centering
%   \caption{Beschreibung der Tabelle}
%   \label{tab:beispiel}
%   \begin{tabular}{lcc}
%     \toprule
%     \textbf{Kategorie} & \textbf{Wert A} & \textbf{Wert B} \\
%     \midrule
%     Zeile 1 & 42 & 58 \\
%     Zeile 2 & 37 & 63 \\
%     \bottomrule
%   \end{tabular}
% \end{table}


% ============================================================
\chapter{Ergebnisse}
\label{chap:ergebnisse}
% ============================================================

% Hinweis: Dieses Kapitel praesentiert die Ergebnisse der Analyse.
% Die Ergebnisse werden strukturiert dargestellt und interpretiert.

[Einfuehrung: Was sind die zentralen Ergebnisse der Untersuchung?]


\section{Darstellung der Ergebnisse}
\label{sec:darstellung}

[Systematische Praesentation der Ergebnisse. Jedes Teilergebnis
wird klar benannt und mit Belegen untermauert.]


\section{Interpretation und Diskussion}
\label{sec:diskussion}

[Einordnung der Ergebnisse in den theoretischen Kontext.
Was bedeuten die Ergebnisse? Wie verhalten sie sich zum
Stand der Forschung? Gibt es ueberraschende Befunde?]

[Kritische Reflexion der eigenen Vorgehensweise und moegliche
Einschraenkungen der Ergebnisse.]

%% ============================================================
%% Fazit
%% ============================================================

\chapter{Fazit}
\label{chap:fazit}

% Hinweis: Das Fazit umfasst typischerweise 2-3 Seiten und enthaelt
% eine Zusammenfassung, die Beantwortung der Forschungsfrage und
% einen Ausblick. Es werden KEINE neuen Informationen eingefuehrt.


\section{Zusammenfassung}
\label{sec:zusammenfassung}

[Kompakte Zusammenfassung der wichtigsten Inhalte und Erkenntnisse
der Arbeit. Was wurde untersucht? Welche Methode wurde angewendet?
Was sind die zentralen Ergebnisse?]

[Die Zusammenfassung soll den roten Faden der Arbeit nachzeichnen,
ohne Details zu wiederholen.]


\section{Beantwortung der Forschungsfrage}
\label{sec:beantwortung}

Die in der Einleitung formulierte Forschungsfrage lautete:

\begin{quote}
\textit{[Forschungsfrage hier nochmals anfuehren.]}
\end{quote}

[Explizite und begruendete Beantwortung der Forschungsfrage auf
Grundlage der erarbeiteten Ergebnisse. Konnte die Frage vollstaendig
beantwortet werden? Welche Teilaspekte konnten geklaert werden?]


\section{Ausblick}
\label{sec:ausblick}

[Welche weitergehenden Fragestellungen ergeben sich aus der Arbeit?
Wo besteht weiterer Forschungsbedarf? Welche Aspekte konnten im
Rahmen dieser VWA nicht behandelt werden und waeren fuer zukuenftige
Untersuchungen interessant?]

[Optional: Persoenliche Reflexion ueber den Arbeitsprozess und
den Erkenntnisgewinn.]


% ============================================================
% === Backmatter ===
% ============================================================

% Literaturverzeichnis
\clearpage
\addcontentsline{toc}{chapter}{Literaturverzeichnis}
\printbibliography[title={Literaturverzeichnis}]

% Abbildungsverzeichnis
\clearpage
\addcontentsline{toc}{chapter}{Abbildungsverzeichnis}
\listoffigures

% Abkuerzungsverzeichnis (falls Eintraege vorhanden)
\printglossary[type=\acronymtype, title={Abk\"urzungsverzeichnis}]

% Glossar (falls Eintraege vorhanden)
\printglossary[title={Glossar}]

% Anhang
\appendix
%% ============================================================
%% Anhang
%% ============================================================

\chapter{Anhang}
\label{chap:anhang}

% Hinweis: Der Anhang enthaelt ergaenzendes Material, das im
% Haupttext den Lesefluss stoeren wuerde, aber fuer die
% Nachvollziehbarkeit der Arbeit wichtig ist.


\section{Interviewleitfaden}
\label{sec:interviewleitfaden}

% Falls Interviews durchgefuehrt wurden, hier den Leitfaden anfuegen.
% Beispiel:

% Das folgende Leitfadeninterview wurde am [Datum] mit [Person/Funktion]
% durchgefuehrt. Die Fragen gliederten sich in folgende Themenbloecke:
%
% \begin{enumerate}
%   \item Einstiegsfragen
%   \begin{itemize}
%     \item Koennten Sie sich kurz vorstellen?
%     \item Wie sind Sie zu diesem Thema gekommen?
%   \end{itemize}
%   \item Hauptfragen
%   \begin{itemize}
%     \item [Frage 1]
%     \item [Frage 2]
%   \end{itemize}
%   \item Abschlussfragen
%   \begin{itemize}
%     \item Gibt es etwas, das Sie noch ergaenzen moechten?
%   \end{itemize}
% \end{enumerate}


\section{Interviewtranskripte}
\label{sec:transkripte}

% Transkripte koennen hier eingefuegt oder als separate Dateien
% eingebunden werden:
% \input{kapitel/transkript-interview1}

[Hier die transkribierten Interviews einfuegen, falls vorhanden.]


\section{Erklaerung zum Einsatz von KI-Werkzeugen}
\label{sec:ki-erklaerung}

% Hinweis: Gemaess den aktuellen Richtlinien des BMBWF muss die
% Verwendung von KI-gestuetzten Werkzeugen (z.B. ChatGPT, DeepL,
% Grammarly) transparent dokumentiert werden.

Im Rahmen dieser Vorwissenschaftlichen Arbeit wurden folgende
KI-gestuetzte Werkzeuge eingesetzt:

\begin{itemize}
  \item \textbf{[Name des Werkzeugs, \zB ChatGPT 4o]:}
    [Beschreibung des Einsatzzwecks, \zB \glqq Wurde zur Ideenfindung
    und Strukturierung von Argumenten verwendet. Alle generierten
    Inhalte wurden kritisch geprueft und eigenstaendig ueberarbeitet.\grqq]
  \item \textbf{[Name des Werkzeugs, \zB DeepL]:}
    [Beschreibung des Einsatzzwecks, \zB \glqq Wurde zur Uebersetzung
    englischsprachiger Fachliteratur als Verstaendnishilfe genutzt.\grqq]
\end{itemize}

% Falls keine KI-Werkzeuge eingesetzt wurden:
% Im Rahmen dieser Vorwissenschaftlichen Arbeit wurden keine
% KI-gestuetzten Werkzeuge eingesetzt.


\section{Weitere Materialien}
\label{sec:weitere-materialien}

% Hier koennen weitere Materialien eingefuegt werden, z.B.:
% - Fragebogen
% - Statistische Auswertungen
% - Abbildungen in hoeherer Aufloesung
% - Quellenabdrucke

[Weitere ergaenzende Materialien hier einfuegen.]


\end{document}
